\documentclass[12pt]{article}
\usepackage[T1]{fontenc}
\usepackage[utf8]{inputenc}
\usepackage{lmodern}
\usepackage{textcomp}
\usepackage{enumitem}
\usepackage{geometry}
\geometry{margin=1in}
\begin{document}
\begin{center}
Answer Key - History - Graduate Level Assessment 2 \
\small (Answers are ordered exactly as in the question paper.)
\end{center}

\section*{Multiple Choice}
\begin{enumerate}[label=\arabic*.]
\item \textbf{Answer:} E
\\ \textit{Options:} A) 1200 cc; B) under 650 cc; C) 1700 cc; D) 1350 cc; E) just under 1000 cc
\\ \textit{Note:} Homo erectus is characterized by a cranial capacity that typically ranges from about 600 to 1,100 cubic centimeters, with the mean being just under 1,000 cc, reflecting its evolutionary advancements in brain size compared to earlier hominins.
\item \textbf{Answer:} A
\\ \textit{Options:} A) Mi'kmaq; B) Norse; C) Beothuk; D) Dorset; E) Nootka
\\ \textit{Note:} Genetic evidence indicates that the modern Inuit share significant ancestral ties with the Mi'kmaq, illustrating a historical lineage that connects these groups through migration and adaptation in Arctic environments.
\item \textbf{Answer:} B
\\ \textit{Options:} A) 5,000,000; 100\%; B) 3,000,000; 75\%; C) 4,000,000; 80\%; D) 2,500,000; 90\%; E) 1,500,000; 70\%
\\ \textit{Note:} The correct answer reflects the historical context of ancient Egypt, where the population was approximately 3 million people, and a significant majority, around 75\%, were engaged in agriculture, highlighting the agrarian nature of the society under the pharaoh's rule.
\item \textbf{Answer:} A
\\ \textit{Options:} A) mutilation and violence.; B) stars and comets.; C) infants and children.; D) modern technologies.; E) domesticated animals.
\\ \textit{Note:} Mayan art and iconography prominently feature depictions of mutilation and violence as they often reflect the civilization's religious beliefs and societal practices, including human sacrifice and the portrayal of warfare, which were integral to their worldview and cultural identity.
\item \textbf{Answer:} C
\\ \textit{Options:} A) the development of language and writing.; B) slowly increased brain size and the first stone tools.; C) a rapid increase in brain size.; D) the transition from bipedalism to quadrupedalism.; E) a decrease in stature and increase in brain size.
\\ \textit{Note:} The correct answer is supported by the fact that after 400,000 years ago, significant changes in hominid brain size occurred, reflecting advancements in cognitive abilities, social structures, and tool use, which are crucial for understanding human evolution.
\end{enumerate}

\section*{True / False}
\begin{enumerate}[label=\arabic*.]
\setcounter{enumi}{5}
\item \textbf{Answer:} False
\\ \textit{Options:} A) True; B) False
\\ \textit{Note:} The correct answer is: from the low-angle perspective of someone swimming
\item \textbf{Answer:} True
\\ \textit{Options:} A) True; B) False
\\ \textit{Note:} According to the passage: sulfonamides
\item \textbf{Answer:} True
\\ \textit{Options:} A) True; B) False
\\ \textit{Note:} According to the passage: light gun or SHORAD
\item \textbf{Answer:} False
\\ \textit{Options:} A) True; B) False
\\ \textit{Note:} The correct answer is: Middle Ages
\item \textbf{Answer:} False
\\ \textit{Options:} A) True; B) False
\\ \textit{Note:} The correct answer is: Roman morality
\end{enumerate}

\section*{Long Form Response}
\begin{enumerate}[label=\arabic*.]
\setcounter{enumi}{10}
\item \textbf{Answer:} Mary Beth Cahill
\\ \textit{Note:} Expected answer: Mary Beth Cahill
\item \textbf{Answer:} Beautiful Stranger
\\ \textit{Note:} Expected answer: Beautiful Stranger
\end{enumerate}

\vfill
\noindent\textit{End of Answer Key}
\end{document}